% chktex-file 13
% chktex-file 36
Báo cáo đã hoàn thành mục tiêu nghiên cứu và triển khai cơ chế hỗ trợ đối tượng \texttt{std::atomic} trong công cụ kiểm thử đơn vị tự động AKAUTAUTO. Thông qua quá trình khảo sát và phân tích kiến trúc hệ thống, nghiên cứu đã chỉ ra các hạn chế của phiên bản hiện tại khi xử lý biến nguyên tử, đặc biệt tại các giai đoạn nhận diện kiểu dữ liệu, biểu diễn dữ liệu kiểm thử và sinh mã kiểm thử.

Để khắc phục những hạn chế này, nghiên cứu đã đề xuất và cài đặt lớp \texttt{AtomicDataNode} nhằm biểu diễn đối tượng \texttt{std::atomic} trong cây dữ liệu kiểm thử. Bên cạnh đó, cơ chế sinh test driver cũng được mở rộng để hỗ trợ khởi tạo và xử lý dữ liệu nguyên tử thông qua các thao tác \texttt{load()} và \texttt{store()}, thay vì sử dụng phép gán trực tiếp như đối với các kiểu dữ liệu thông thường. Giải pháp được tích hợp vào kiến trúc hiện có của AKAUTAUTO mà không làm thay đổi các thành phần cốt lõi của hệ thống.

Kết quả thực nghiệm trên bộ chương trình gồm 25 hàm sử dụng \texttt{std::atomic} dưới nhiều hình thức khác nhau cho thấy hệ thống có thể nhận diện thành công các đối tượng nguyên tử, xây dựng chính xác \texttt{AtomicDataNode}, sinh mã kiểm thử hợp lệ và thực thi thành công các ca kiểm thử tương ứng. Các trường hợp thực nghiệm bao gồm tham số truyền tham chiếu, con trỏ, bí danh kiểu dữ liệu, đối tượng nguyên tử trong \texttt{struct}/\texttt{class}, các thao tác nguyên tử phổ biến như \texttt{load()}, \texttt{store()}, \texttt{exchange()}, \texttt{fetch\_*} và \texttt{compare\_exchange()}, cũng như các kịch bản liên quan đến cơ chế stub. Những kết quả này cho thấy giải pháp đề xuất đã mở rộng đáng kể phạm vi hỗ trợ của AKAUTAUTO đối với các chương trình C++ có sử dụng đối tượng nguyên tử.

Trong tương lai, nghiên cứu có thể được mở rộng theo nhiều hướng khác nhau. Thứ nhất, hệ thống có thể tiếp tục hỗ trợ đầy đủ hơn các thao tác của thư viện \texttt{std::atomic}, bao gồm các hàm nguyên tử nâng cao và các kiểu dữ liệu nguyên tử phức tạp hơn. Thứ hai, cơ chế sinh dữ liệu kiểm thử tự động có thể được tích hợp với bộ giải ràng buộc nhằm tự động tìm kiếm các giá trị đầu vào giúp kích hoạt những nhánh điều kiện phụ thuộc vào trạng thái của đối tượng nguyên tử. Cuối cùng, việc hỗ trợ các kịch bản đa luồng và phân tích tương tác đồng thời giữa nhiều đối tượng \texttt{atomic} là một hướng phát triển tiềm năng, góp phần nâng cao khả năng kiểm thử của AKAUTAUTO đối với các hệ thống phần mềm hiện đại.
